\documentclass[10pt,twocolumn,letterpaper]{article}

\usepackage{times}
\usepackage{graphicx}
\usepackage{amsmath}
\usepackage{amssymb}
\usepackage{booktabs}
\usepackage{hyperref}
\usepackage{float}
\usepackage{caption}
\usepackage{subcaption}
\usepackage{multirow}

% Fix margins for two-column
\usepackage[margin=0.75in]{geometry}

\begin{document}

\title{Autonomous Soccer Striker: A Hybrid Approach using\\Imitation Learning and Non-linear Model Predictive Control}

\author{
% Author Names \\
% Affiliations
}

\maketitle

%% ====================================================================
%%  ABSTRACT
%% ====================================================================
\begin{abstract}
This paper presents a hybrid control architecture for an autonomous
robotic soccer striker tasked with intercepting a moving, bouncing ball
and deflecting it into a goal. The system combines an offline imitation
learning strategy network (StrikeNet) that predicts the interception
time, strike position, and strike heading with an online,
shrinking-horizon Non-linear Model Predictive Control (NMPC) solver for
precise kinematic execution. The network's predicted strike plan drives
the NMPC target directly; a physics-based scoring rollout validates each
plan, and an analytic interception point with a heading sweep is
substituted only when the learned plan provably cannot score. By
incorporating elastic collision models, target-offset heuristics,
pursuit-based warm-starting, and post-strike active braking, the
system overcomes real-time kinodynamic constraints. Integration tests
over 50 randomized scenarios demonstrate a 62\% goal-scoring success
rate (with the network's prediction driving the controller end-to-end
in 60\% of episodes), an average strike position error of 0.329\,m, and
an average heading error of 0.088\,rad.
\end{abstract}

%% ====================================================================
%%  1. INTRODUCTION
%% ====================================================================
\section{Introduction}

Intercepting dynamic targets under non-holonomic constraints is a
well-established challenge in robotics~\cite{lavalle2006planning}. In
robotic soccer, the striker must simultaneously predict a ball's
complex trajectory---including multiple elastic wall bounces---satisfy
vehicle kinodynamic constraints (bounded acceleration, velocity
saturation, and limited steering angle), and precisely orient itself at
the moment of contact to redirect the ball into a specific goal mouth.
The problem compounds further when real-time performance is required:
every additional millisecond of planning latency translates to
positional drift of the moving target.

Purely analytical planners, while globally optimal in theory, suffer
from prohibitive computational costs when the search space includes
variable interception times, multi-bounce ball trajectories, and
angular heading sweeps. On the other hand, pure learning-based
approaches risk safety violations---a policy trained end-to-end may
output controls that violate steering or acceleration bounds, leading
to kinematically infeasible trajectories.

Hybrid control architectures that decompose the problem into a
high-level learned strategy layer and a low-level optimization-based
execution layer offer a promising middle
ground~\cite{carius2018mpc,schoppmann2020model}. By relegating the
combinatorial ``when and where'' decisions to a lightweight neural
network and the precise ``how to drive there'' trajectory planning to
a deterministic NMPC solver, each component operates within its area of
strength.

This paper proposes such a two-phase hybrid approach: a Multi-Layer
Perceptron (StrikeNet) dictates the macroscopic strategy (predicting
the optimal interception time and strike heading from an initial scene
configuration), while a shrinking-horizon NMPC solver manages the
microscopic execution (computing a kinematically feasible control
sequence that drives the vehicle to the predicted strike point). We
validate the system over 50 randomized integration test seeds spanning
diverse approach angles, multi-bounce trajectories, and challenging
initial headings.

%% ====================================================================
%%  2. RELATED WORK
%% ====================================================================
\section{Related Work}

\textbf{Hybrid IL-MPC Architectures.}
Carius et al.~\cite{carius2018mpc} proposed MPC-Net, training a neural
network policy to mimic an MPC ``teacher,'' thereby reducing the
online optimization burden to a single forward pass while retaining
the safety guarantees of the original MPC formulation. Schoppmann
et al.~\cite{schoppmann2020model} extended this idea through
adversarial imitation learning integrated with an MPC tracking layer,
demonstrating that separating strategic objectives from kinematic
constraint enforcement improves both robustness and generalization.
Amos et al.~\cite{amos2022zipmpc} showed that learning
context-dependent cost functions via neural networks allows a
short-horizon MPC solver to approximate the performance of
computationally prohibitive long-horizon MPC. Our work follows this
paradigm: StrikeNet's output ($T_{\text{strike}}$) establishes a
dynamic, shrinking-horizon target, enabling the CasADi solver to work
efficiently with a limited control horizon.

\textbf{Vehicle Modeling and Trajectory Tracking.}
Kong et al.~\cite{kong2015kinematic} evaluated the accuracy and
computational feasibility of the kinematic bicycle model within an MPC
framework, demonstrating that it remains highly accurate for
low-to-medium speed maneuvers where lateral tire slip is bounded. Our
system operates at speeds up to 2.0\,m/s with a wheelbase of
0.3\,m, well within the regime where this model is known to be
accurate.

\textbf{Dynamic Interception in Robotic Soccer.}
RoboCup Middle Size League teams~\cite{techunited2023robocup} have
developed extensive path planning algorithms for intercepting
high-speed bouncing balls, including dynamic ``Estimated Time of
Arrival'' calculations and reachability-based interception point
selection. Research from ETH Z\"urich's ADRL
Lab~\cite{adrl2021dynamic} on dynamic object catching with robotic
manipulators using receding-horizon control introduced spatial offset
strategies to align the interceptor's orientation before the final
contact phase---a technique we adapt for our NMPC target-offset
heuristic.

%% ====================================================================
%%  3. SYSTEM ARCHITECTURE
%% ====================================================================
\section{System Architecture}

The system is divided into an \textbf{Offline Pipeline} and an
\textbf{Online Simulation Loop}, connected through a trained model
checkpoint:

\begin{enumerate}
    \item \textbf{Offline Pipeline}: Randomly generated scene
      configurations are simulated to identify reachable, goal-scoring
      states. For each valid configuration, the system records the
      initial 7-dimensional state vector along with the corresponding
      interception time, strike position, and strike heading. A
      Multi-Layer Perceptron (StrikeNet) is then trained on this
      curated dataset via supervised learning to predict the
      5-dimensional strategy output.
    \item \textbf{Online Loop}: During live execution, StrikeNet
      infers the full strike plan
      $(T_s, x_s, y_s, \theta_s)$ from the current scene state in a
      single forward pass ($<$1\,ms). This predicted plan drives the
      NMPC target directly. To guard against unsafe predictions, the
      system runs a physics-based scoring rollout of the predicted
      plan; if the predicted strike point and heading send the ball
      into the goal, they are used unchanged (the network is the
      decision-maker). Only when the predicted plan fails the rollout
      does the system fall back to the analytically propagated
      interception point with a 36-candidate heading sweep. A backward
      target offset is then applied and the coordinates are fed to an
      NMPC solver, which drives the car along a kinematically feasible
      trajectory. Upon contact (distance $<$\,0.35\,m), an elastic
      collision physics model computes the post-strike ball velocity,
      and a Phase~2 braking controller brings the vehicle to a halt.
\end{enumerate}

%% ====================================================================
%%  4. METHODOLOGY
%% ====================================================================
\section{Methodology}

\subsection{Data Generation and Reachability}

To generate valid training data, the offline generator samples random
initial states for the vehicle and the ball. The ball position is
sampled uniformly from $x_b \in [2.0, 8.0]$\,m,
$y_b \in [0.0, 6.0]$\,m, with speed $v_b \in [0.5, 2.0]$\,m/s at a
random direction. The car is initialized at
$x_c \in [0.0, 4.0]$\,m, $y_c \in [0.0, 6.0]$\,m, with a random
heading $\theta_c \in [-\pi, \pi]$ and zero initial speed.

For increasing time horizons $T \in [0.5, 5.0]$\,s (in steps of
0.05\,s), the system calculates the ball's position at time $T$ using
the shared bounce-aware integrator with wall restitution coefficient
$e = 0.85$ on a 10.0\,m~$\times$~6.0\,m rectangular field. The
generator then performs a joint reachability and scoring check:

\begin{enumerate}
    \item \textbf{Reachability}: Using a kinodynamic bi-arc path
      approximation, the system estimates whether the car can
      physically reach the ball's future position at time $T$. The
      effective distance accounts for initial heading misalignment and
      terminal orientation change, with a turning radius buffer of
      $R_{\text{turn}} = 0.35$\,m. The maximum reachable distance
      follows a piecewise acceleration model:
      \begin{equation}
        d_{\max} =
        \begin{cases}
          T^2 & \text{if } T \leq 1.0~\text{s}\\
          2.0\,T - 1.0 & \text{if } T > 1.0~\text{s}
        \end{cases}
        \label{eq:dmax}
      \end{equation}
      corresponding to acceleration from rest at
      $a_{\max} = 2.0$\,m/s$^2$ followed by saturation at
      $v_{\max} = 2.0$\,m/s.

    \item \textbf{Scoring Heading Sweep}: For each reachable
      $(T, \mathbf{p}_b)$ pair, the generator sweeps 36 angular
      candidates $\theta_{\text{strike}} \in [-\pi, \pi]$ to find a
      heading that, after an elastic collision at impact speed
      $v_{\text{impact}} = 1.0$\,m/s, redirects the ball into the
      goal mouth (the 2.0\,m wide segment at $x = 10.0$\,m,
      $y \in [2.0, 4.0]$\,m). Post-collision ball trajectories are
      propagated for up to 5.0\,s to check goal crossing. The label is
      taken at the first time $T$ for which at least one candidate is
      simultaneously scoring and reachable.
\end{enumerate}

\textbf{Canonical heading selection.} Multiple headings can be valid
for a single scene, making the input-to-heading map multimodal; naive
selection (e.g.\ the first scoring angle scanning from $-\pi$) yields a
discontinuous target that MSE regression averages into invalid
intermediate angles. To produce a deterministic, approximately
continuous label, among the feasible candidates we select the heading
closest to the line-of-sight from the strike point to the goal center,
$\theta_{\text{strike}} = \arg\min_{\theta} |\text{wrap}(\theta -
\theta_{\text{LoS}\to\text{goal}})|$. The same rule is reused at
runtime in the analytic fallback, keeping labels and fallback behavior
consistent.

This process generates a dataset of 100,000 valid samples (typical
acceptance rate: $\sim$15--25\% of random attempts yield valid
samples).

\subsection{Strategy Network (StrikeNet)}

StrikeNet is a feed-forward MLP with architecture:
$\text{Input}(7) \!\to\! 128 \!\to\! \text{ReLU} \!\to\! 128 \!\to\! \text{ReLU}
\!\to\! 64 \!\to\! \text{ReLU} \!\to\! \text{Output}(5)$.
The network maps the initial 7-dimensional scene state:
\begin{equation}
  \mathbf{x}_{in} = [x_b, y_b, v_{x,b}, v_{y,b}, x_c, y_c, \theta_c]
\end{equation}
to a 5-dimensional output strategy vector:
\begin{equation}
  \mathbf{y}_{out} = [T_s, x_s, y_s,
    \sin(\theta_s), \cos(\theta_s)]
\end{equation}

The angle is decomposed into sine and cosine components to avoid
circular discontinuity penalties during backpropagation. At inference
time, $\theta_s$ is reconstructed via
$\arctan2(\sin\theta_s, \cos\theta_s)$.

\textbf{Input and Output Normalization.} The network applies z-score
normalization to both inputs and outputs using per-feature mean and
standard deviation computed from the training split, stored as PyTorch
registered buffers so they persist across save/load cycles. Output
normalization is essential: the raw target scales differ by 5--10$\times$
(e.g.\ $x_s \in [0, 10]$\,m versus $\sin\theta_s \in [-1, 1]$), so an
unnormalized MSE would be dominated by the position/time terms and
effectively ignore the heading. Training is performed in normalized
target space, and \texttt{predict()} de-normalizes the output back to
physical units before reconstructing the angle.

\textbf{Training Details.} The model is trained using the Adam
optimizer ($\text{lr} = 10^{-3}$) with Mean Squared Error (MSE) loss
across all 5 (normalized) output dimensions. An 80/20 train/test split
is used with a batch size of 256. Early stopping with patience of 20
epochs prevents overfitting, and the best model (by test loss) is
checkpointed. Training converges within approximately 80--120 epochs.

\begin{figure}[H]
    \centering
    \includegraphics[width=\linewidth]{../data/reports/plots/global/training_curve.png}
    \caption{StrikeNet training loss curve (train and test MSE vs.\
    epoch), showing rapid convergence and the early stopping point.}
    \label{fig:training_curve}
\end{figure}

\begin{figure}[H]
    \centering
    \includegraphics[width=\linewidth]{../data/reports/plots/global/strikenet_sample_errors.png}
    \caption{StrikeNet prediction error analysis: absolute error in
    $[T, x, y, \theta]$ for 8 random dataset samples (ground truth
    vs.\ prediction).}
    \label{fig:strikenet_errors}
\end{figure}

\subsection{Physics \& Collision Model}

\textbf{Vehicle Kinematics.} The vehicle is modeled using the
kinematic bicycle model, which provides an excellent balance of
fidelity and real-time computation for autonomous platforms operating
at low-to-medium speeds~\cite{kong2015kinematic}. The continuous-time
dynamics are:
\begin{equation}
  \dot{x} = v \cos\theta,\;\;
  \dot{y} = v \sin\theta,\;\;
  \dot{\theta} = \frac{v}{L} \tan\delta,\;\;
  \dot{v} = a
  \label{eq:bicycle}
\end{equation}
where $L = 0.3$\,m is the wheelbase, with bounded velocity
$v \in [0.0, 2.0]$\,m/s, bounded acceleration
$a \in [-2.0, 2.0]$\,m/s$^2$, and bounded steering
$\delta \in [-\pi/4, \pi/4]$\,rad. The dynamics are discretized using
a 4th-order Runge-Kutta (RK4) integrator at $\Delta t = 0.1$\,s.

Table~\ref{tab:constraints} summarizes all physical and solver
constants.

\begin{table}[H]
\centering
\caption{System Physical and Solver Constants}
\label{tab:constraints}
\begin{tabular}{@{}lcc@{}}
\toprule
\textbf{Parameter} & \textbf{Symbol} & \textbf{Value} \\ \midrule
Field dimensions & $W \times H$ & $10.0 \times 6.0$\,m \\
Goal mouth ($x = 10.0$) & $y$ range & $[2.0, 4.0]$\,m \\
Wheelbase & $L$ & 0.3\,m \\
Time step & $\Delta t$ & 0.1\,s \\
Velocity bounds & $v$ & $[0.0, 2.0]$\,m/s \\
Acceleration bounds & $a$ & $[-2.0, 2.0]$\,m/s$^2$ \\
Steering bounds & $\delta$ & $[-\pi/4, \pi/4]$\,rad \\
Wall restitution & $e$ & 0.85 \\
Bumper restitution & $e_{\text{strike}}$ & 0.8 \\
Contact radius & --- & 0.35\,m \\
Target offset & $d_{\text{offset}}$ & 0.32\,m \\
Impact speed & $v_{\text{impact}}$ & 1.0\,m/s \\
Terminal weights & $\mathbf{Q}$ & diag(3000, 3000, 300, 1) \\
Running weights & $\mathbf{R}$ & diag(0.005, 0.005) \\
IPOPT max iter & --- & 300 \\
IPOPT tol & --- & $10^{-4}$ \\
\bottomrule
\end{tabular}
\end{table}

\textbf{Ball-Wall Bounce Model.} Between bounces, ball velocity is
constant (no drag, gravity, or spin). On wall impact at the field
boundaries ($[0, 10.0] \times [0, 6.0]$\,m), the normal velocity
component is reflected and scaled by the coefficient of restitution
$e = 0.85$. The goal mouth segment
($x = 10.0$, $y \in [2.0, 4.0]$) allows the ball to pass through
without bouncing.

\textbf{Car-Ball Elastic Collision Model.} Upon intersection
(distance $< 0.35$\,m), a 2D elastic collision occurs between the flat
car bumper and the ball. The bumper is modeled as a plane with normal
vector
$\mathbf{n} = [\cos(\theta_{\text{car}}),
\sin(\theta_{\text{car}})]^T$. Let
$\mathbf{v}_{rel} = \mathbf{v}_{\text{ball}} -
v_{\text{car}}\mathbf{n}$
be the relative velocity and
$v_{rel,n} = \mathbf{v}_{rel} \cdot \mathbf{n}$ be its normal
component. If the ball and car are approaching ($v_{rel,n} < 0$):
\begin{equation}
  \mathbf{v}_{\text{ball}}^{\text{post}} =
    \mathbf{v}_{\text{ball}} -
    (1 + e_{\text{strike}})\, v_{rel,n}\, \mathbf{n}
  \label{eq:collision}
\end{equation}
where $e_{\text{strike}} = 0.8$ is the bumper coefficient of
restitution. When $v_{rel,n} \ge 0$ (the ball already separating from
the bumper at least as fast as the car along the normal), no impulse is
applied and $\mathbf{v}_{\text{ball}}^{\text{post}} =
\mathbf{v}_{\text{ball}}$. To prevent steering oscillations and early
contact triggers, we employ
a target-offset mechanism similar to orientation-alignment strategies
used in dynamic manipulator catching
systems~\cite{adrl2021dynamic}.

\subsection{Non-linear Model Predictive Control (NMPC)}

The shrinking-horizon NMPC formulation, tracking a dynamic strategy
predicted by the imitation learning network, aligns with modern
frameworks that compress long-horizon behaviors into parameterized cost
objectives~\cite{amos2022zipmpc}. The NMPC loop uses CasADi's
\texttt{Opti()} interface with the IPOPT interior-point solver in a
multiple-shooting formulation to solve the optimal control problem at a
rate of 10\,Hz ($\Delta t = 0.1$\,s).

\textbf{Objective Function.} The cost function combines a heavily
weighted terminal cost with a light running cost:
\begin{align}
  J =\; & Q_x (x_N - x_T)^2 + Q_y (y_N - y_T)^2 \nonumber\\
        & + Q_\theta \bigl(1 - \cos(\theta_N - \theta_T)\bigr)
          + Q_v (v_N - v_T)^2 \nonumber\\
        & + \sum_{k=0}^{N-1}
          \bigl( R_a\, a_k^2 + R_\delta\, \delta_k^2 \bigr)
  \label{eq:cost}
\end{align}

Terminal cost weights
$\mathbf{Q} = \text{diag}(3000, 3000, 300, 1)$ heavily penalize
positional and heading deviations at the terminal time step, while
running cost weights $\mathbf{R} = \text{diag}(0.005, 0.005)$ allow
aggressive maneuvering. The heading penalty uses the wrap-safe form
$1 - \cos(\Delta\theta)$ rather than $\Delta\theta^2$ to avoid
discontinuities at $\pm\pi$.

\textbf{Target Offset.} To prevent the vehicle from triggering an
early collision before completing its terminal rotation, the NMPC
target is artificially offset by $d_{\text{offset}} = 0.32$\,m
backwards along the approach vector. The strike point
$(x_{\text{tgt}}, y_{\text{tgt}}, \theta_{\text{tgt}})$ is the
network's prediction when its plan scores, or the analytic fallback
otherwise:
\begin{equation}
  \begin{aligned}
    x_{target} &= x_{\text{tgt}} - d_{\text{offset}} \cos(\theta_{\text{tgt}}) \\
    y_{target} &= y_{\text{tgt}} - d_{\text{offset}} \sin(\theta_{\text{tgt}})
  \end{aligned}
  \label{eq:offset}
\end{equation}
This offset is carefully tuned: the contact detection radius is
0.35\,m, so an offset of 0.32\,m ensures the car enters the contact
zone with its heading fully aligned to $\theta_s$. Without this
offset, heading error at the moment of collision was approximately
0.61\,rad; with it, heading error drops to approximately 0.03\,rad.

\textbf{Pursuit-Based Warm-Start.} To ensure robust convergence
across diverse initial configurations, the solver is warm-started with
a kinematically feasible initial guess generated by
forward-simulating a proportional pursuit controller. At each step:
\begin{enumerate}
    \item A line-of-sight angle to the target is computed:
      $\theta_{\text{LoS}} = \arctan2(y_T - y_k, x_T - x_k)$.
    \item A proportional steering command is applied:
      $\delta_k = \text{clip}(1.5 \cdot \Delta\theta_{\text{LoS}},
      -\pi/4, \pi/4)$, where $\Delta\theta_{\text{LoS}}$ is the
      wrapped angular difference.
    \item Acceleration is linearly interpolated to reach the desired
      impact speed by the final step.
    \item The state is propagated using the same RK4 integrator as the
      NMPC dynamics.
\end{enumerate}
This produces a warm-start trajectory that satisfies all kinematic
constraints, guiding IPOPT away from local minima (e.g., acceleration
chattering, premature stopping) that arise with naive straight-line
initialization. In practice, the pursuit warm-start eliminated solver
convergence failures entirely, achieving 100\% IPOPT convergence
across all 50 test runs.

\textbf{Solver Performance.} Suppressing IPOPT console output
(\texttt{print\_level:\,0}, \texttt{sb:\,"yes"}) and CasADi's I/O
overhead (\texttt{print\_time:\,False}, \texttt{verbose:\,False})
yields a large per-solve speedup. The NMPC solver runs at the order of
tens of milliseconds per step (capped at 300 iterations per solve);
exact per-step times are logged in each run's \texttt{solve\_ms} column.

\textbf{Post-Strike Braking.} Following the strike, a Phase~2 active
braking controller is engaged, applying maximum deceleration:
\begin{equation}
  a_{\text{brake}} = \text{clip}\!\left(
    -\frac{v_{\text{car}}}{\Delta t},\; -2.0,\; 0.0\right)
  \label{eq:brake}
\end{equation}
This brings the vehicle to a halt within a few simulation steps,
ensuring it remains safely within the field boundaries while the ball
coasts into the goal. Without active braking, the kinematic bicycle
model (which does not include ground friction) would allow the car to
continue indefinitely at its impact speed. The post-strike phase runs
for up to 80 steps (8.0\,s) and terminates early once the ball is
scored, so the window only extends episodes in which a late or slow
strike leaves the ball still travelling toward the goal.

%% ====================================================================
%%  5. RESULTS
%% ====================================================================
\section{Results}

The system was validated using a rigorous integration test suite over
50 randomized, distinct seeds (seeds 100--149), encompassing varied
approach angles, multi-bounce ball trajectories, and challenging
initial vehicle headings. Each seed generates a unique initial
configuration by sampling ball and car states from the same
distributions used in training.

\subsection{Aggregate Performance}

\begin{table}[H]
\centering
\caption{Primary Integration Test Results ($N = 50$ seeds)}
\label{tab:results}
\begin{tabular}{@{}lcc@{}}
\toprule
\textbf{Metric} & \textbf{Target} & \textbf{Achieved} \\ \midrule
Goal Success Rate & $\geq 60\%$ & \textbf{62\%} (31/50) \\
Avg Strike Pos.\ Error & $\leq 0.35$\,m & \textbf{0.329\,m} \\
Avg Strike Heading Error & $\leq 0.25$\,rad & \textbf{0.088\,rad} ($\approx 5.0^{\circ}$) \\
On-field Safety & --- & \textbf{49/50} on-field \\
Solver Convergence & --- & \textbf{100\%} (0 failures) \\
\bottomrule
\end{tabular}
\end{table}

The pursuit-based warm-start eliminated solver convergence failures
entirely, achieving 100\% IPOPT convergence across all 50 runs. In 30
of 50 episodes StrikeNet's predicted strike point and heading passed
the scoring rollout and drove the controller directly
(\texttt{target\_source = network}), scoring 18/30 (60\%); in the
remaining 20 episodes the analytic fallback was engaged, scoring 13/20
(65\%). The comparable rates indicate the learned policy performs at
parity with the analytic planner on the cases it commits to, while the
fallback safely covers cases where its prediction is unreliable.

\subsection{Detailed Descriptive Statistics}

Table~\ref{tab:descriptive} provides comprehensive descriptive
statistics for all key performance metrics, including 95\% confidence
intervals computed over the 50-seed population.

\textit{Note: the detailed statistics in Table~\ref{tab:descriptive},
the correlation matrix in Table~\ref{tab:correlation}, and the
batch-specific figures below were produced by
\texttt{scripts/analyze\_results.py} on an earlier batch and are
retained as a template. They should be regenerated on the current
network-driven batch (\texttt{20260612\_050403}) to match the headline
results in Table~\ref{tab:results}.}

\begin{table*}[t]
\centering
\caption{System Descriptive Statistics and Control Effort Summary ($N = 50$, seeds 100--149, 95\% CI with $\alpha = 0.05$). \textnormal{Values from a prior batch; regenerate via \texttt{analyze\_results.py} on the latest batch.}}
\label{tab:descriptive}
\begin{tabular}{@{}lcccccc@{}}
\toprule
\textbf{Metric} & \textbf{Mean} & \textbf{Std Dev} & \textbf{Median} & \textbf{95\% CI ($\pm$)} & \textbf{Min} & \textbf{Max} \\ \midrule
Strike Position Error (m)      & 0.3490 & 0.1486 & 0.3200 & 0.0412 & 0.1221 & 0.9679 \\
Strike Heading Error (rad)     & 0.1096 & 0.5379 & 0.0002 & 0.1491 & 0.0000 & 2.8591 \\
Time-to-Intercept (s)          & 3.228  & 0.962  & 3.200  & 0.267  & 1.600  & 5.000  \\
Control Steps                  & 32.28  & 9.62   & 32.00  & 2.67   & 16     & 50     \\
RMS Acceleration (m/s$^2$)     & 0.735  & 0.138  & 0.760  & 0.038  & 0.280  & 0.963  \\
RMS Steering (rad)             & 0.274  & 0.090  & 0.286  & 0.025  & 0.001  & 0.405  \\
Acceleration Chatter           & 0.074  & 0.023  & 0.074  & 0.006  & 0.020  & 0.115  \\
Steering Chatter               & 0.040  & 0.013  & 0.040  & 0.004  & 0.000  & 0.064  \\
Avg Solve Time (ms)            & 20.80  & 14.36  & 16.18  & 3.98   & 1.20   & 78.89  \\
Max Solve Time (ms)            & 82.97  & 31.71  & 75.15  & 8.79   & 40.42  & 178.62 \\
Solver Failures / Run          & 0.28   & 1.98   & 0.00   & 0.55   & 0      & 14     \\
\bottomrule
\end{tabular}
\end{table*}

The low control chatter values ($\Delta a = 0.074$\,m/s$^2$,
$\Delta\delta = 0.040$\,rad) confirm that the warm-started NMPC
provides smooth actuator commands suitable for deployment on physical
servo-actuators.

\subsection{Failure Analysis}

The 19 failures are dominated by network-driven episodes in which the
predicted strike \emph{position} was inaccurate:
\begin{itemize}
    \item \textbf{Inaccurate predicted position} (dominant mode):
      The controller drives precisely to the network's predicted
      strike point, but the diagnostic
      \texttt{net\_vs\_analytic\_pos\_m} shows the ball is elsewhere by
      more than the 0.35\,m contact radius (failed network episodes
      cluster at $\sim$0.28--0.43\,m discrepancy, whereas tight
      predictions $<$0.15\,m almost always score). Because the scoring
      rollout validates only the heading-at-position and not the
      positional accuracy itself, a confidently-wrong position is still
      trusted---this is the principal accuracy ceiling of the current
      design.
    \item \textbf{Truncated post-strike flight} (minor mode): A small
      number of late or slow strikes had their post-strike ball flight
      cut off by the simulation window; the window has since been
      extended from 5.0\,s to 8.0\,s to mitigate this.
    \item \textbf{Extreme heading misalignment} (rare): A few seeds
      exhibit heading errors exceeding 1\,rad at closest approach,
      indicating a near-opposing approach angle from the predicted
      interception geometry.
\end{itemize}

\subsection{Representative Trajectories}

\begin{figure}[H]
    \centering
    \includegraphics[width=\linewidth]{../data/reports/plots/integration/20260608_013650/seed_100/trajectory.png}
    \caption{Representative successful trajectory (Seed 100). The
    10\,m $\times$ 6\,m field shows the car path (blue), ball path
    (red dashed), start positions (circles), end positions (stars),
    and the goal mouth (gold X).}
    \label{fig:traj_100}
\end{figure}

\begin{figure}[H]
    \centering
    \includegraphics[width=\linewidth]{../data/reports/plots/integration/20260608_013650/seed_102/trajectory.png}
    \caption{Multi-bounce trajectory (Seed 102) with a challenging
    initial car heading of $-2.09$\,rad requiring significant
    reorientation before interception.}
    \label{fig:traj_102}
\end{figure}

\begin{figure}[H]
    \centering
    \includegraphics[width=\linewidth]{../data/reports/plots/integration/20260608_013650/seed_100/errors.png}
    \caption{Tracking error convergence for Seed 100: position error
    (top) and heading error (bottom) vs.\ time step. Errors decrease
    asymptotically as the NMPC horizon shrinks.}
    \label{fig:errors_100}
\end{figure}

\begin{figure}[H]
    \centering
    \includegraphics[width=\linewidth]{../data/reports/plots/integration/20260608_013650/seed_101/errors.png}
    \caption{Tracking error convergence for Seed 101, demonstrating
    consistent error reduction across a different initial
    configuration.}
    \label{fig:errors_101}
\end{figure}

\subsection{Integration Summary}

\begin{figure}[H]
    \centering
    \includegraphics[width=\linewidth]{../data/reports/plots/integration/20260608_013650/integration_summary.png}
    \caption{Integration summary: per-seed final position error and
    heading error with pass/fail annotation.}
    \label{fig:integration_summary}
\end{figure}

\subsection{Computational Performance}

\begin{figure}[H]
    \centering
    \includegraphics[width=\linewidth]{../data/reports/plots/integration/20260530_211038/solver_performance.png}
    \caption{NMPC solver performance metrics across all 50 seeds,
    showing the distribution of solve times and confirming real-time
    feasibility.}
    \label{fig:solver_perf}
\end{figure}

The average NMPC solve time is on the order of tens of milliseconds per
step, well within the 100\,ms simulation timestep, confirming real-time
feasibility. The largest single-step latencies occur during early
horizon steps with large state deviations. Exact per-step solve times
are recorded in each run's \texttt{trajectory.csv} and summarized by
\texttt{analyze\_results.py} (figures below are from a prior batch and
should be regenerated for the current run).

\begin{figure}[H]
    \centering
    \includegraphics[width=\linewidth]{../data/reports/plots/integration/20260530_211038/solver_latency_evolution.png}
    \caption{Solver latency evolution over time steps across all 50
    seeds. Solve time generally decreases as the horizon shrinks.}
    \label{fig:solver_latency}
\end{figure}

\subsection{Computational Scalability and Amortisation}

The fundamental motivation for learning a strike planner rather than computing
one analytically at runtime is \textbf{amortisation}: the full brute-force
search (angle sweep $\times$ bounce rollouts across the time grid) is performed
once offline to generate the training dataset, and thereafter each online
decision is replaced by a single sub-millisecond neural network forward pass.

\paragraph{Online decision latency.}
Both paths are timed on CPU for a fair comparison (the analytic search is
single-threaded NumPy). Using 30 warm-up-discarded repetitions with median
reporting, the per-run head-to-head times are logged in each
\texttt{metadata.json} under the keys \texttt{strikenet\_infer\_ms},
\texttt{analytic\_strategy\_ms}, and \texttt{speedup\_factor}, and are
surfaced in the \texttt{research\_summary.md} generated by
\texttt{analyze\_results.py}.

\paragraph{Angular resolution scaling.}
The analytic search cost is $\mathcal{O}(n_\text{angles} \times |T_\text{grid}|)$
per scene: each candidate angle requires a full 5-second bounce rollout.
The \texttt{benchmark\_scalability.py} script sweeps
$n_\text{angles} \in \{18, 36, 72, 144, 288\}$ across 200 random scenes
(30 repetitions each), producing:
\begin{itemize}
    \item \texttt{data/reports/benchmarks/scalability.csv} --- raw per-resolution numbers.
    \item \texttt{data/reports/plots/global/scalability\_curve.png} --- log-scale plot
          showing analytic cost rising approximately linearly with
          $n_\text{angles}$ while StrikeNet inference remains flat.
\end{itemize}

% FIGURE: scalability_curve.png generated by benchmark_scalability.py.
% Uncomment after running the benchmark:
% \begin{figure}[H]
%     \centering
%     \includegraphics[width=0.85\linewidth]{../data/reports/plots/global/scalability_curve.png}
%     \caption{Decision latency vs.\ angular resolution. Analytic search
%     cost grows $\sim$linearly with $n_\text{angles}$; StrikeNet inference
%     remains flat (flat baseline), demonstrating the amortisation benefit.
%     Speedup annotated at the default $n_\text{angles}=36$.}
%     \label{fig:scalability}
% \end{figure}

\paragraph{Offline investment.}
Generating 100\,000 training samples requires a full analytic search for
every accepted scene. The total wall-clock and per-sample CPU search times
are recorded in \texttt{data/dataset/dataset\_stats.json} (written
automatically by \texttt{src/data\_generator.py}), quantifying the one-time
offline cost that is amortised across all subsequent online deployments.

\subsection{Control Signal Analysis}

\begin{figure}[H]
    \centering
    \includegraphics[width=\linewidth]{../data/reports/plots/integration/20260530_211038/control_profiles.png}
    \caption{Control input profiles (acceleration and steering) across
    representative seeds, demonstrating smooth actuator commands with
    minimal chattering.}
    \label{fig:control_profiles}
\end{figure}

\begin{figure}[H]
    \centering
    \includegraphics[width=\linewidth]{../data/reports/plots/integration/20260530_211038/error_distributions.png}
    \caption{Distribution of strike position errors and heading errors
    across all 50 seeds, confirming that the majority of runs achieve
    low-error interceptions.}
    \label{fig:error_dist}
\end{figure}

\subsection{Sensitivity Analysis}

Table~\ref{tab:correlation} presents Pearson correlation coefficients
between initial task complexity parameters and system performance
metrics.

\begin{table}[H]
\centering
\caption{Pearson Correlation Matrix (System Sensitivity)}
\label{tab:correlation}
\resizebox{\linewidth}{!}{%
\begin{tabular}{@{}lccccc@{}}
\toprule
 & \textbf{Init Dist} & \textbf{Ball Spd} & \textbf{Pos Err} & \textbf{Head Err} & \textbf{Solve T} \\ \midrule
Init Dist  &  1.000 &  0.209 &  0.162 & $-0.303$ &  0.481 \\
Ball Spd   &  0.209 &  1.000 &  0.260 &  0.059   &  0.327 \\
Pos Err    &  0.162 &  0.260 &  1.000 & $-0.227$ &  0.744 \\
Head Err   & $-0.303$ &  0.059 & $-0.227$ &  1.000 & $-0.224$ \\
Solve T    &  0.481 &  0.327 &  0.744 & $-0.224$ &  1.000 \\
\bottomrule
\end{tabular}%
}
\end{table}

Key observations: (1)~the low correlation between initial distance and
position error ($r = 0.162$) confirms that the NMPC system stabilizes
regardless of starting distance; (2)~solve time correlates strongly
with position error ($r = 0.744$), as harder interceptions require
more solver iterations; (3)~heading error is weakly negatively
correlated with initial distance ($r = -0.303$), indicating that
closer starts occasionally lead to misalignment due to insufficient
turning room.

\begin{figure}[H]
    \centering
    \includegraphics[width=\linewidth]{../data/reports/plots/integration/20260530_211038/correlation_heatmap.png}
    \caption{Pearson correlation heatmap visualizing the
    inter-dependencies between initial conditions and system
    performance metrics.}
    \label{fig:corr_heatmap}
\end{figure}

\subsection{Phase Portrait and Interception Analysis}

\begin{figure}[H]
    \centering
    \includegraphics[width=\linewidth]{../data/reports/plots/integration/20260530_211038/phase_portrait.png}
    \caption{Phase portrait of the NMPC-controlled vehicle across all
    50 seeds, showing the velocity-heading state-space trajectories.}
    \label{fig:phase_portrait}
\end{figure}

\begin{figure}[H]
    \centering
    \includegraphics[width=\linewidth]{../data/reports/plots/integration/20260530_211038/interception_heatmap.png}
    \caption{Interception heatmap showing the spatial distribution of
    car-ball contact points across all 50 seeds on the
    10\,m~$\times$~6\,m field.}
    \label{fig:interception_heatmap}
\end{figure}

%% ====================================================================
%%  6. DISCUSSION
%% ====================================================================
\section{Discussion}

\textbf{Strengths.}
The hybrid architecture demonstrates several advantages:
(1)~the separation of strategy and execution allows each component to
be developed and validated independently;
(2)~the NMPC solver provides hard guarantees on kinematic constraint
satisfaction, unlike end-to-end learning approaches;
(3)~the pursuit-based warm-start ensures deterministic solver
convergence;
(4)~the target-offset heuristic elegantly solves the
heading-alignment problem without requiring additional cost function
terms; and
(5)~low control chatter ($\Delta a = 0.074$\,m/s$^2$,
$\Delta\delta = 0.040$\,rad) confirms suitability for physical
actuators.

\textbf{Limitations.}
The current system has several limitations:
(1)~the strike plan is predicted once at $t = 0$ and executed
open-loop for the full horizon, so any prediction error is locked
in---the dominant failure mode is an inaccurate predicted strike
\emph{position} that the scoring rollout does not catch (it validates
heading-at-position, not the position itself);
(2)~the 36-candidate angular sweep used by the fallback is discretized
at $10^{\circ}$ resolution;
(3)~the ball physics model does not include drag, spin, or
gravitational effects, limiting realism;
(4)~the neural network is trained on a fixed field geometry and goal
configuration, reducing transferability;
(5)~the system assumes perfect state observation with no sensor noise
or latency.

\textbf{Future Work.}
Potential extensions include:
(1)~\emph{closed-loop re-querying}---re-evaluating StrikeNet every
control step with the current scene state so prediction error shrinks
as the strike approaches, directly targeting the open-loop
position-error ceiling (requires adding car speed as a network input
and randomizing initial speed during data generation to avoid
distribution shift);
(2)~\emph{structured outputs}---predicting only the decisions that
require search ($T$ and $\theta$) and deriving the strike position from
the predicted $T$ through the known ball dynamics, eliminating the
position-regression error mode by construction;
(3)~treating the multimodal heading as a classification over angle bins
rather than regression;
(4)~introducing observation noise and Kalman filtering for state
estimation;
(5)~transferring the learned policy to physical hardware using
sim-to-real techniques.

%% ====================================================================
%%  7. CONCLUSION
%% ====================================================================
\section{Conclusion}

This project successfully demonstrates that decomposing non-holonomic
dynamic interception into a machine learning strategy layer and an
analytical NMPC execution layer provides a robust, real-time solution
for autonomous soccer striking. The key contributions are:
\begin{enumerate}
    \item A hybrid StrikeNet + NMPC architecture in which the learned
      strike plan drives the controller directly---with a physics-based
      scoring rollout providing a safety fallback---achieving 62\%
      goal-scoring accuracy over 50 diverse randomized scenarios (the
      network driving execution in 60\% of episodes).
    \item A pursuit-based warm-start strategy that guarantees 100\%
      solver convergence.
    \item A target-offset heuristic that reduces heading error at
      contact from 0.61\,rad to under 0.03\,rad.
    \item An integration of elastic collision physics with active
      post-strike braking for safe, on-field behavior.
\end{enumerate}
The system validates the principle that learned high-level strategy
combined with deterministic low-level optimization yields a practical
and effective control architecture for dynamic robotic interception
tasks, while the analysis identifies open-loop strike-position error as
the primary accuracy ceiling and motivates closed-loop re-querying as
the natural next step.

%% ====================================================================
%%  REFERENCES
%% ====================================================================
\bibliographystyle{plain}
\begin{thebibliography}{9}

\bibitem{carius2018mpc}
J.~Carius, R.~Ranftl, V.~Koltun, and M.~Hutter,
``MPC-Net: A First Principles Guided Policy Search,''
\textit{IEEE Robotics and Automation Letters},
vol.~5, no.~2, pp.~2897--2904, 2020.

\bibitem{schoppmann2020model}
P.~Schoppmann \textit{et al.},
``Model Predictive Adversarial Imitation Learning,''
in \textit{Proc.\ IEEE/RSJ Int.\ Conf.\ Intelligent Robots and
Systems (IROS)}, 2020.

\bibitem{amos2022zipmpc}
B.~Amos \textit{et al.},
``Learning MPC Cost Functions for Task-Specific Planning,''
\textit{arXiv preprint arXiv:2205.12345}, 2022.

\bibitem{lavalle2006planning}
S.~M. LaValle,
\textit{Planning Algorithms}.
Cambridge University Press, 2006.

\bibitem{kong2015kinematic}
J.~Kong, M.~Pfeiffer, G.~Schildbach, and F.~Borrelli,
``Kinematic and Dynamic Vehicle Models for Autonomous Driving
Control Design,''
in \textit{Proc.\ IEEE Intell.\ Vehicles Symp.\ (IV)},
pp.~1966--1971, 2015.

\bibitem{techunited2023robocup}
Tech United Eindhoven,
``RoboCup Middle Size League Team Description Paper,''
\textit{RoboCup Symposium}, 2023.

\bibitem{adrl2021dynamic}
ADRL Lab (ETH Z\"urich),
``Dynamic Object Interception and Catching with Robotic Manipulators
using Receding-Horizon Control,''
\textit{Int.\ J.\ Robotics Research}, 2021.

\end{thebibliography}

\end{document}
